% Техническое_задание.tex — Техническое задание (DOC-02), версия 0.5
% Комплект документации криптотренажёра mephi-12 (НИЯУ МИФИ, курсовая работа)
% Стиль — ДО-178C/КТ-178C; преамбула и титул — по единому шаблону full_doc_template.tex
% Термины — дословно из сквозного глоссария (методика идентификаторов v1.0);
% факты о Системе — из карточек доменов и каталога эндпоинтов v0.1 (Спринт 1).
\documentclass[12pt,a4paper]{article}

% =============================================================================
% КОНФИГУРАЦИЯ И ПРЕАМБУЛА
% =============================================================================

% --- конфигурация метаданных ---
% ---------- Прямые значения ----------
\newcommand{\Company}{НИЯУ МИФИ}
\newcommand{\DocumentTitleMain}{Техническое задание}
\newcommand{\DocumentTitleSub}{на разработку платформы для изучения и практики курса по криптографии в МИФИ}
\newcommand{\Version}{0.5}
\newcommand{\DocDate}{24.10.2026}
\newcommand{\AuthorName}{Кокорев В.Д.}
\newcommand{\AgreedBy}{}
\newcommand{\ApprovedBy}{}
\newcommand{\City}{Москва}

% --- преамбула и стили ---
\usepackage[utf8]{inputenc}
\usepackage[T2A]{fontenc}
\usepackage[russian]{babel}

% Times New Roman-like
\usepackage{newtxtext,newtxmath} % похож на Times; обеспечивает текст+математику
\usepackage{setspace}
\onehalfspacing % межстрочный интервал 1.5
\setlength{\emergencystretch}{2em} % допуск для строк с длинными кодовыми фрагментами

% Страницы
\usepackage{geometry}
\geometry{left=25mm,right=15mm,top=20mm,bottom=20mm}

% отступ первой строки
\usepackage{indentfirst}
\setlength{\parindent}{1.25cm}

% таблицы, картинки, списки, подписи
\usepackage{graphicx}
\usepackage{caption}
\usepackage{subcaption}
\captionsetup[table]{position=top}    % подпись таблицы над ней
\captionsetup[figure]{position=bottom}% подпись рисунка под ней
\usepackage{booktabs}
\usepackage{tabularx}
\usepackage{longtable}
\usepackage{array}
\usepackage{enumitem}
\setlist[itemize]{left=0pt,label=\textbullet}
\setlist[enumerate]{left=0pt,label=\arabic*.}

% код и конфигурации
\usepackage{listings}
\lstset{basicstyle=\ttfamily\small,breaklines=true}

% гиперссылки
\usepackage{hyperref}
\hypersetup{colorlinks=true,linkcolor=black,urlcolor=blue,
  pdftitle={Техническое задание — криптотренажёр mephi-12},
  pdfauthor={Кокорев В.Д.}}

\pagestyle{plain}  % только номер страницы внизу по центру

% нумерация разделов — стандартная десятичная (article по умолчанию)
% дополнительные удобства
\usepackage{tocloft} % управление оглавлением
\setlength{\cftbeforesecskip}{6pt}

% небольшие утилиты
\newcommand{\reqfield}[1]{\textbf{#1}:} % для шаблонов в приложениях

\begin{document}

% =============================================================================
% ТИТУЛЬНЫЙ ЛИСТ
% =============================================================================
\begin{titlepage}
  \thispagestyle{empty} % убирает номер/колонтитул на этой странице
  \centering
  \vspace*{5.8cm}
  {\Huge \textbf{\DocumentTitleMain}\par}
  {\Large \vspace{4pt} \DocumentTitleSub \par} % вторая часть — обычным шрифтом, чуть ниже
  \vspace{0.8cm}

  {\large Версия: \Version \par}
  {\large Дата: \DocDate \par}
  \vspace{1.2cm}

  % ----- Таблица с ролями -----
  \centering
  \begin{tabularx}{0.95\textwidth}{@{\hspace{6pt}}p{5.5cm} p{6cm} p{4cm}@{\hspace{6pt}}}
    \toprule
    \textbf{Роль} & \textbf{ФИО} & \textbf{Подпись / Дата} \\
    \midrule
    Разработал: & \AuthorName & \\
    \addlinespace[6pt]
    Согласовано: & \AgreedBy & \\
    \addlinespace[6pt]
    Утверждено: & \ApprovedBy & \\
    \bottomrule
  \end{tabularx}

  \vfill
  {\large г. \City, \the\year \par}
\end{titlepage}

% =============================================================================
% ЛИСТ РЕГИСТРАЦИИ ИЗМЕНЕНИЙ
% =============================================================================
\section*{Лист регистрации изменений \texttt{[TS.CHANGES]}}
\addcontentsline{toc}{section}{Лист регистрации изменений}

\begin{table}[h!]
\centering
\small
\begin{tabularx}{\textwidth}{
  >{\centering\arraybackslash}p{1.1cm}
  >{\centering\arraybackslash}p{2.1cm}
  p{2.3cm}
  p{2.6cm}
  p{1.6cm}
  X
  p{1.8cm}
}
\toprule
\textbf{Версия} & \textbf{Дата} & \textbf{Автор} & \textbf{Изменённый элемент} & \textbf{Х-р изменения} & \textbf{Описание изменения} & \textbf{Утвердил} \\
\midrule
0.5 & 24.10.2026 & Кокорев В.Д. & Весь документ & Добавлено & Первичный выпуск (черновая версия Спринта 1). Детализация требований — Спринт 2; v1.0 — до 30.10.2026. & \\
\bottomrule
\end{tabularx}
\end{table}

\clearpage
\tableofcontents
\clearpage

% =============================================================================
% ОСНОВНОЕ СОДЕРЖАНИЕ
% =============================================================================

\section{Введение \texttt{[TS.INTRO]}}

\subsection{Наименование \texttt{[TS.INTRO.NAME]}}
\begin{description}[style=nextline]
    \item[Полное наименование:] «Платформа для изучения и практики курса по криптографии в МИФИ».
    \item[Краткое наименование:] криптотренажёр mephi-12.
\end{description}
Далее по тексту используется наименование «Система» (определение — в п.~1.4).

\subsection{Цели и назначение разработки \texttt{[TS.INTRO.PURPOSE]}}
Система предназначена для практической подготовки студентов НИЯУ МИФИ по криптографическим методам защиты информации — ранцевым криптосистемам (на супервозрастающей последовательности и на последовательности степеней), бэкдор-атаке на ранцевую криптосистему и схеме Эль-Гамаля. Подготовка осуществляется через решение автоматически генерируемых и автоматически проверяемых на сервере задач, объединяемых в сессии задач с ограничением времени, в контексте учебных групп и именованных конфигураций наборов задач.

\paragraph{Цели разработки.}
\begin{itemize}
    \item Автоматическая генерация учебных криптографических задач и автоматическая проверка решений на сервере, без участия преподавателя в рутинной проверке.
    \item Организация практики в форме сессий задач с лимитом времени на задачу и автоматическим провалом просроченных задач.
    \item Учёт результатов решения по учебным группам средствами конфигураций наборов задач и функции модерации (просмотр сессий студентов).
    \item Разграничение возможностей пользователей по ролям (USER, ADMIN).
    \item Демонстрация на практике уязвимостей криптосистем — встроенного бэкдора в ранцевой криптосистеме — как элемента компетенций по криптоанализу.
\end{itemize}

\subsection{Область применения \texttt{[TS.INTRO.SCOPE]}}
Система применяется в учебном процессе кафедры НИЯУ МИФИ: поддержка лабораторных работ и практикума по дисциплинам, включающим криптографические методы защиты информации. Система эксплуатируется как серверное приложение (backend) учебного назначения; взаимодействие с ней осуществляется по REST API (порт 8083, базовый путь \texttt{/api}); клиентская часть в объём настоящего технического задания не входит.

\subsection{Термины и определения \texttt{[TS.INTRO.TERMS]}}
Применяемые в настоящем техническом задании термины соответствуют сквозному глоссарию комплекта документации (методика идентификаторов v1.0); ниже приведены основные термины; полные канонические определения — в глоссарии (\texttt{notes/methodology/id-methodology.md}).

\begin{description}[style=nextline]
    \item[Система:] «Платформа для изучения и практики курса по криптографии в МИФИ» (учебный криптотренажёр, репозиторий mephi-12/backend): монолит Kotlin 1.9.25 + Spring Boot 3.4.0 (Web, JPA, Security, Validation) с БД PostgreSQL 17, развёртываемый средствами Docker/docker-compose; REST API — базовый путь \texttt{/api}, порт 8083.
    \item[Пользователь:] Учётная запись (сущность \texttt{User}): уникальный email, имя, BCrypt-хеш пароля; опционально привязан к одной группе (\texttt{User.group}); идентифицируется в запросах по клейму \texttt{userId} своего JWT.
    \item[Студент:] Пользователь в роли USER, решающий сессии задач; все операции создания сессий и отправки решений выполняются от имени студента (владельца активного JWT).
    \item[Роль:] Значение authority в JWT: \texttt{ROLE\_USER} (выдаётся каждому при регистрации/логине) и \texttt{ADMIN}. \texttt{SecurityConfig}: пути \texttt{/api/auth/**} и \texttt{/api/public/**} открыты, все остальные требуют роль USER.
    \item[JWT:] Access-токен стандарта JSON Web Token (библиотека jjwt), выпускается при регистрации и логине (\texttt{JwtHelper}); клеймы: \texttt{userId}, \texttt{authorities}, \texttt{permissions}; срок жизни задаётся переменной окружения \texttt{ACCESS\_TOKEN\_LIFETIME} (в днях); проверяется фильтром \texttt{JwtFilter} на каждом запросе.
    \item[Сессия задач:] Период решения студентом набора задач (сущность \texttt{ProblemSession}): создаётся \texttt{POST /api/tasks/\allowbreak session?configuration=<имя>} по включённой для группы студента конфигурации; у студента не более одной активной сессии (\texttt{sessionState = NEW}, повторное создание отклоняется); время жизни задачи — 20 минут, по истечении задача автоматически проваливается планировщиком автопровала; сессия закрывается, когда не осталось задач в состоянии NEW.
    \item[Задача:] Элемент сессии (сущность \texttt{Problem}): JSON-условие (\texttt{statement}), тип (\texttt{ProblemType}) и состояние (\texttt{ProblemState}: NEW «Задание решается», SOLVED «Задание верно решено», FAILED «Задание решено неверно»); проверка решения — \texttt{PUT /api/tasks/\allowbreak submit?statementId=<UUID>}; владелец и активность сессии проверяются при отправке.
    \item[Тип задачи:] Значение перечисления \texttt{ProblemType}: \texttt{BACKPACK\_CODE\_\allowbreak SUPER\_INCREASING} (ранцевая КС на супервозрастающей последовательности), \texttt{BACKPACK\_CODE\_\allowbreak DEGREES} (ранцевая КС на степенях), \texttt{BACKPACK} (ранцевая КС в составе сессии), \texttt{BACKPACK\_BACK\allowbreak DOOR} (бэкдор в ранцевой КС), \texttt{EL\_GAMAL} (КС Эль-Гамаля). В конфигурациях сессий задачи задаются строковыми квалификаторами «Backpack» и «ElGamal»; типы \texttt{BACKPACK\_CODE\_\allowbreak *} используются редакционными (editorial) генерациями \texttt{GET/POST /api/tasks/\allowbreak backpack/editorial}.
    \item[Ранцевая криптосистема (КС):] Криптосистема Меркля--Хеллмана на задаче о рюкзаке: лёгкий рюкзак $L$ (супервозрастающая последовательность), модуль $n = \Sigma L + 1$, множитель $\omega$, взаимно простой с $n$; публичный тяжёлый рюкзак $H = (L \cdot \omega) \bmod n$; шифрование — сумма элементов $H$ по единичным битам сообщения; расшифрование — умножением на $\omega^{-1} \bmod n$.
    \item[Супервозрастающая последовательность:] Последовательность, каждый элемент которой строго больше суммы всех предыдущих; в типе \texttt{BACKPACK\_CODE\_\allowbreak SUPER\_INCREASING} из неё строится лёгкий рюкзак.
    \item[Бэкдор:] Секрет, встроенный в публичный ключ ранцевой КС на степенях (эндпоинт \texttt{GET /api/backdoor/editorial}): десятичные записи $\omega$ и $n$ «зашиты» в младшие разряды элементов тяжёлого рюкзака (цифра $H_i \bmod 10$), что позволяет владельцу секрета восстановить закрытый ключ из открытого.
    \item[Схема Эль-Гамаля:] Криптосистема на мультипликативной группе $Z_p^{*}$: ключи ($p$, $g$, закрытый $x$, открытый $y = g^x \bmod p$), шифрование с сеансовым показателем $k$, расшифрование через $x$; учебные генерации — \texttt{GET /api/tasks/\allowbreak eigamal?bitLength=<N>} (задача) и \texttt{GET /api/tasks/\allowbreak eigamal/editorial?bitLength=<N>} (разбор).
    \item[Конфигурация задач:] Сущность \texttt{ProblemConfiguration}: имя + упорядоченный список квалификаторов типов задач (\texttt{taskServices}), задающий состав и порядок задач в создаваемой сессии; доступность группе — через \texttt{UserGroupProblem\allowbreak Configuration.enabled} (включённые конфигурации группы: \texttt{GET /api/tasks/groups/\allowbreak configurations?groupId=<UUID>}).
    \item[Группа пользователей:] Сущность \texttt{UserGroup} (учебная группа): объединяет пользователей (\texttt{User.group}) и связывается с конфигурациями задач; список групп — \texttt{GET /api/users/groups}, привязка себя к группе — \texttt{PUT /api/users/\allowbreak me/groups?group=<имя>}.
    \item[Планировщик автопровала:] Фоновый процесс \texttt{ProblemAutoFailScheduler} (раз в 60~с): переводит задачи в состоянии NEW старше 20 минут (по времени создания) в FAILED и закрывает сессии, в которых не осталось задач в состоянии NEW.
\end{description}

\section{Общее описание системы \texttt{[TS.DESC]}}

\subsection{Концепция \texttt{[TS.DESC.CONCEPT]}}
Система реализуется как монолитное серверное приложение со слоистой архитектурой: REST-сервис на Kotlin 1.9.25 / Spring Boot 3.4.0 (стартеры Web, JPA, Security, Validation) с СУБД PostgreSQL 17; развёртывание — средствами Docker/docker-compose (контейнер БД с контролем готовности, затем контейнер приложения). Система предоставляет 14 действующих эндпоинтов в 7 контроллерах.

Студенты регистрируются и вступают в учебные группы. С каждой группой связываются именованные конфигурации задач — упорядоченные наборы типов задач; студент запускает сессию задач по включённой конфигурации своей группы и получает автоматически сгенерированные условия. Решения отправляются на сервер и проверяются автоматически, с вынесением вердикта по каждой задаче (SOLVED — верно решено, FAILED — решено неверно). На решение каждой задачи отводится 20 минут; по истечении лимита задача автоматически проваливается фоновым планировщиком, что исключает неограниченное время решения.

Роли: USER — все операции студента (решение задач); ADMIN — административные функции, MODERATOR как отдельная роль не выделяется (функция модерации — просмотр сессий всех студентов); об ограничениях текущей реализации доступа — п.~2.3.

\subsection{Архитектура и взаимодействие \texttt{[TS.DESC.ARCH]}}

\paragraph{Основной сценарий использования.}
\begin{enumerate}
    \item \textbf{Регистрация:} студент передаёт email, имя и пароль; сервер валидирует поля (email 4--120 символов, пароль 8--255) и при успехе сразу выдаёт JWT.
    \item \textbf{Логин:} при повторном входе студент аутентифицируется парой «email + пароль» и получает JWT.
    \item \textbf{Аутентифицированный доступ:} все дальнейшие запросы передаются с заголовком \texttt{Authorization:\allowbreak\ Bearer <JWT>}; клеймы \texttt{userId} и \texttt{authorities} идентифицируют студента и его роль.
    \item \textbf{Выбор группы:} студент просматривает список учебных групп и вступает в одну из них по названию; без группы сессия не создаётся.
    \item \textbf{Просмотр конфигураций:} студент просматривает конфигурации задач своей группы (имя набора, состав типов, признак включения).
    \item \textbf{Запуск сессии:} студент создаёт сессию по имени конфигурации; сервер генерирует задачи в порядке, заданном конфигурацией (квалификаторы «Backpack» и «ElGamal»), и возвращает сессию в состоянии NEW со списком задач.
    \item \textbf{Получение задачи:} условие каждой задачи (JSON), тип и состояние NEW входят в состав сессии; тип задачи определяется конфигурацией (ранцевая КС либо схема Эль-Гамаля).
    \item \textbf{Решение:} студент готовит решение — восстанавливает или вычисляет параметры криптосистемы по условию.
    \item \textbf{Отправка решения:} решение отправляется по идентификатору задачи (\texttt{statementId}); сервер проверяет владельца и активность сессии и передаёт решение проверщику соответствующего типа.
    \item \textbf{Вердикт:} задача переводится в SOLVED либо в FAILED (с диагностическим сообщением о неверном параметре); когда задач в состоянии NEW не остаётся, сессия завершается (исчерпание задач либо автопровал по таймеру — см. фоновые сценарии).
\end{enumerate}

\paragraph{Типы задач (что решает студент).}
\begin{description}[style=nextline]
    \item[\texttt{BACKPACK\_CODE\_SUPER\_INCREASING} (ранцевая КС на супервозрастающей последовательности):] студент получает условие с частично скрытым лёгким ранцем (часть элементов скрыта, часть параметров отсутствует) и восстанавливает все параметры криптосистемы: лёгкий ранец, модуль, множитель $\omega$, тяжёлый ранец, обратный множитель $\omega^{-1}$, шифртекст и расшифрованное сообщение (расшифрование — жадным алгоритмом по супервозрастающей последовательности).
    \item[\texttt{BACKPACK\_CODE\_DEGREES} (ранцевая КС на последовательности степеней):] то же, но лёгкий ранец строится из элементов вида $p^i$ (степени основания $p$); студент выполняет шифрование и расшифрование в этой конструкции, восстанавливая тот же набор параметров.
    \item[\texttt{BACKPACK\_BACKDOOR} (бэкдор в ранцевой КС, учебная форма атаки Шамира):] студент получает публичный тяжёлый ранец, в младшие разряды элементов которого «зашиты» десятичные записи секретных $\omega$ и $n$, и убеждается, что владелец секрета восстанавливает закрытый ключ из открытого. Тип используется только редакционной демонстрацией и в сессиях не участвует.
    \item[\texttt{EL\_GAMAL} (схема Эль-Гамаля):] студент получает открытые параметры и сообщение ($p$, $g$, $x$, $m$), вычисляет открытый ключ $y = g^x \bmod p$, выбирает сеансовый показатель $k$, строит компоненты шифртекста $c_1$ и $c_2$ и выполняет расшифрование; проверка решения — по шагам, с сообщением о первом неверном промежуточном значении.
\end{description}

\paragraph{Фоновые сценарии.}
Планировщик автопровала выполняется каждые 60~с: задачи в состоянии NEW, созданные более 20 минут назад, переводятся в FAILED (момент истечения определяется по времени создания задачи); сессии, в которых не осталось задач в состоянии NEW, закрываются. Автопровал не зависит от действий студента, в том числе от его отключения от Системы.

\paragraph{Учебные (editorial) сценарии.}
Редакционные эндпоинты предназначены для самоподготовки и подготовки методических материалов: эталонная ранцевая задача выбранного типа с полным решением и возможностью задания параметров (\texttt{GET/POST /api/tasks/\allowbreak backpack/editorial}), демонстрационный пример со встроенным бэкдором (\texttt{GET /api/backdoor/editorial}), пошаговый разбор схемы Эль-Гамаля (\texttt{GET /api/tasks/\allowbreak eigamal/\allowbreak editorial}) и задание Эль-Гамаля для самостоятельного решения (\texttt{GET /api/tasks/\allowbreak eigamal}). Результаты редакционных генераций формируются на лету и не записываются в базу данных.

\subsection{Известные ограничения реализации \texttt{[TS.DESC.LIMIT]}}
Ниже честно фиксируются ограничения подтверждённой вычиткой кода реализации; они являются материалом для формальных инспекций и ревизии требований Спринта 2 (детализация — в СТ DOC-05).
\begin{enumerate}
    \item \textbf{Роль ADMIN не назначается:} JWT всегда выпускается без административной authority, т.е. получить роль ADMIN через API невозможно; эндпоинт модерации \texttt{GET /api/moderator/\allowbreak tasks/session} в текущей реализации доступен любому аутентифицированному пользователю с ролью USER.
    \item \textbf{История решений:} \texttt{GET /api/users/me/tasks} обслуживается устаревшим сервисом (помечен как deprecated) и возвращает все сохранённые ранцевые задачи, а не только задачи текущего студента.
    \item \textbf{Управление группами и конфигурациями:} API покрывает только чтение конфигураций и вступление в группу; создание, изменение и удаление групп и конфигураций, а также переключение признака включения выполняются на уровне базы данных.
    \item \textbf{Жёстко зафиксированные параметры:} лимит времени на задачу (20 минут) зашит в SQL-запросе планировщика и не вынесен в конфигурацию; параметры генерации задач Эль-Гамаля в сессии зафиксированы (bitLength = 5) и не настраиваются конфигурацией.
    \item \textbf{Коды состояния при некорректном JWT:} при некорректном JWT (неверная подпись, нарушение формата, отсутствие префикса Bearer) и при недостаточности роли Система возвращает 500 вместо 401/403 (обработчик AppException с состоянием по умолчанию); приведение кодов к семантике REST — кандидат в замечания Спринта 2.
\end{enumerate}

\section{Требования к системе \texttt{[TS.REQ]}}
В настоящей редакции (v0.5) требования приводятся тезисно, на уровне намерений разработки, и соответствуют фактическим возможностям реализованной Системы, кроме расхождений, зафиксированных в п.~2.3 (в частности, по кодам состояния). Полный набор требований с атомарной декомпозицией и присвоением идентификаторов SR-XXX формируется в Спринте 2 в системных требованиях (DOC-05).

\subsection{Функциональные требования \texttt{[TS.REQ.FUNC]}}

\paragraph{\texttt{[TS.REQ.FUNC.AUTH.REG]}.} Система должна обеспечивать регистрацию нового пользователя по адресу электронной почты, имени и паролю с серверной валидацией полей (длина email 4--120 символов, пароля 8--255 символов) и отклонением повторной регистрации по занятому адресу. Успешная регистрация должна завершаться немедленной выдачей JWT.

\paragraph{\texttt{[TS.REQ.FUNC.AUTH.LOGIN]}.} Система должна обеспечивать аутентификацию (логин) по паре «email + пароль» со сверкой пароля по хешу и выдачей JWT; при неверном пароле — отказ с кодом 401.

\paragraph{\texttt{[TS.REQ.FUNC.AUTH.JWT]}.} Система должна требовать предъявления действительного JWT (заголовок \texttt{Authorization: Bearer}) для всех операций, кроме регистрации и логина; клеймы \texttt{userId} и \texttt{authorities} должны идентифицировать текущего студента (профиль пользователя) и его роль во всех доменах Системы.

\paragraph{\texttt{[TS.REQ.FUNC.GROUP.LIST]}.} Система должна предоставлять перечень всех учебных групп (идентификатор, название, состав участников).

\paragraph{\texttt{[TS.REQ.FUNC.GROUP.JOIN]}.} Система должна позволять студенту вступить в учебную группу по её названию (привязка — к одной группе) и должна отклонять создание сессии, если студент не состоит в группе.

\paragraph{\texttt{[TS.REQ.FUNC.CONFIG.VIEW]}.} Система должна предоставлять по группе перечень конфигураций задач (имя, состав типов, признак включения); создание сессии должно допускаться только по включённой конфигурации группы студента.

\paragraph{\texttt{[TS.REQ.FUNC.SESSION.CREATE]}.} Система должна создавать сессию задач по имени конфигурации: генерировать условия задач в порядке, заданном конфигурацией, и возвращать студенту сессию в состоянии NEW с JSON-условиями, типами и состояниями задач. При наличии у студента активной (незавершённой) сессии создание новой должно отклоняться.

\paragraph{\texttt{[TS.REQ.FUNC.SESSION.TIME]}.} Система должна ограничивать время решения каждой задачи 20 минутами и по истечении лимита автоматически переводить задачу в FAILED фоновым планировщиком с периодом срабатывания 60~с, независимо от действий студента.

\paragraph{\texttt{[TS.REQ.FUNC.SESSION.SUBMIT]}.} Система должна принимать решение задачи по её идентификатору (\texttt{statementId}) при условии, что задача принадлежит текущему студенту и его сессия активна, передавать решение проверщику соответствующего типа и выносить вердикт SOLVED либо FAILED с диагностическим сообщением при ошибке. Сессия должна закрываться, когда в ней не остаётся задач в состоянии NEW.

\paragraph{\texttt{[TS.REQ.FUNC.PROB.BACKPACK.SI]}.} Для задач «ранцевая КС на супервозрастающей последовательности» Система должна генерировать условие с частично скрытыми параметрами и проверять решение по всем восстанавливаемым величинам: элементам лёгкого ранца (строго сверхвозрастающим), модулю, взаимной простоте множителя $\omega$ с модулем, тяжёлому ранцу, обратному множителю, шифртексту и расшифрованному сообщению.

\paragraph{\texttt{[TS.REQ.FUNC.PROB.BACKPACK.DEG]}.} Для задач «ранцевая КС на степенях» Система должна строить лёгкий ранец из степеней основания $p$ и выполнять генерацию условия и проверку решения по той же схеме параметров ранцевой КС.

\paragraph{\texttt{[TS.REQ.FUNC.PROB.BACKDOOR]}.} Система должна генерировать демонстрационный пример ранцевой КС со встроенным бэкдором: десятичные записи секретных $\omega$ и модуля размещаются в младших разрядах элементов тяжёлого ранца; при невозможности построить пример за отведённое число попыток — возвращать диагностическую ошибку. Тип используется только редакционной генерацией, в сессиях не участвует.

\paragraph{\texttt{[TS.REQ.FUNC.PROB.ELGAMAL]}.} Для задач «схема Эль-Гамаля» Система должна генерировать учебное задание ($p$, $g$, $x$, $m$) и проверять решение по шагам: открытый ключ $y$, сеансовый показатель $k$, компоненты шифртекста $c_1$ и $c_2$, расшифрованное сообщение; при ошибке — возвращать сообщение о первом неверном промежуточном значении.

\paragraph{\texttt{[TS.REQ.FUNC.EDITORIAL]}.} Система должна предоставлять редакционные (editorial) генерации для самоподготовки: эталонную ранцевую задачу выбранного типа с полным решением (с возможностью задания параметров), демонстрационный пример с бэкдором, пошаговый разбор схемы Эль-Гамаля и задание Эль-Гамаля для самостоятельного решения. Результаты редакционных генераций не должны записываться в базу данных.

\paragraph{\texttt{[TS.REQ.FUNC.MOD.VIEW]}.} Система должна предоставлять функцию модерации — просмотр сессий задач всех студентов (состояния сессий и их задач) для контроля учебного прогресса.

\paragraph{\texttt{[TS.REQ.FUNC.HISTORY]}.} Система должна предоставлять студенту просмотр истории решённых задач (об известном ограничении текущей реализации — п.~2.3).

\subsection{Требования к интерфейсам \texttt{[TS.REQ.INT]}}

\paragraph{\texttt{[TS.REQ.INT.REST]}.} Все внешние интерфейсы Системы должны следовать единой конвенции REST: базовый путь \texttt{/api} (порт 8083), формат обмена данными — JSON, идентификация ресурсов — по UUID.

\paragraph{\texttt{[TS.REQ.INT.BEARER]}.} Аутентификация запросов должна выполняться заголовком \texttt{Authorization: Bearer <JWT>}; без аутентификации должны быть доступны только операции регистрации и логина (\texttt{/api/auth/**}).

\paragraph{\texttt{[TS.REQ.INT.CODES]}.} Система должна единообразно применять коды состояния: 400 — невалидное тело запроса; 401 — отсутствие или истечение JWT, неверные учётные данные; 404 — ресурс не найден. Ошибки должны возвращаться в едином JSON-формате (код, сообщение, время).

\subsection{Требования к надёжности \texttt{[TS.REQ.REL]}}

\paragraph{\texttt{[TS.REQ.REL.AUTOFAIL]}.} Автоматический провал просроченных задач должен обеспечиваться фоновым планировщиком и не должен зависеть от действий студента, в том числе от его отключения от Системы.

\paragraph{\texttt{[TS.REQ.REL.PERSIST]}.} Состояние пользователей, групп, конфигураций, сессий и задач должно храниться в PostgreSQL и сохраняться при перезапуске приложения; после перезапуска незавершённые сессии должны продолжать обслуживаться, включая автопровал по таймеру.

\paragraph{\texttt{[TS.REQ.REL.SCHED]}.} Истечение времени задачи должно определяться по времени её создания, а не по факту срабатывания планировщика: пропуск или задержка отдельного цикла планировщика не должны приводить к некорректным вердиктам — задача проваливается при первом же цикле после истечения лимита.

\paragraph{\texttt{[TS.REQ.REL.DEPLOY]}.} Развёртывание docker-compose должно обеспечивать старт приложения только после готовности СУБД (healthcheck) и сохранность данных базы между перезапусками (том данных).

\subsection{Требования к безопасности \texttt{[TS.REQ.SEC]}}

\paragraph{\texttt{[TS.REQ.SEC.JWT]}.} JWT должен подписываться алгоритмом HMAC-SHA256; срок жизни токена должен задаваться конфигурацией (в днях) через переменную окружения.

\paragraph{\texttt{[TS.REQ.SEC.PASSWORDS]}.} Пароли пользователей не должны храниться и передаваться в открытом виде: хранение — только в виде BCrypt-хеша, сверка — по хешу.

\paragraph{\texttt{[TS.REQ.SEC.ROLES]}.} Система должна разграничивать возможности по ролям (USER, ADMIN), передаваемым в клеймах \texttt{authorities} токена; фактическое назначение роли ADMIN в текущей реализации — см. п.~2.3, уточнение правил доступа — предмет СТ.

\paragraph{\texttt{[TS.REQ.SEC.SECRETS]}.} Секреты и параметры окружения (секрет подписи JWT \texttt{JWT\_SECRET}, срок жизни токена \texttt{ACCESS\_TOKEN\_LIFETIME}, параметры подключения к базе данных) должны передаваться только через переменные окружения и не должны включаться в код или репозиторий.

\subsection{Требования к производительности \texttt{[TS.REQ.PERF]}}

\paragraph{\texttt{[TS.REQ.PERF.CRUD]}.} Операции регистрации, логина, просмотра групп и конфигураций, создания сессии и отправки решения должны выполняться за время, не создающее ощутимых задержек для интерактивной учебной работы.

\paragraph{\texttt{[TS.REQ.PERF.GEN]}.} Генерация учебных задач, включая генерации с параметром bitLength и переборное построение примера с бэкдором (ограниченное число попыток), должна укладываться в интерактивно приемлемое для учебного процесса время.

\textit{Примечание: количественные показатели производительности (время отклика, пропускная способность) в настоящей редакции не устанавливаются; качественные требования уточняются количественными показателями в СТ (DOC-05, раздел \texttt{[SR.PERF]}).}

\section{Требования к документации и разработке \texttt{[TS.DOC]}}

\paragraph{\texttt{[TS.DOC.SET]}.} Разработка должна сопровождаться комплектом документации в стиле ДО-178C/КТ-178C, охватывающим цепочку «требования $\to$ реализация $\to$ верификация»: DOC-00 (README комплекта), DOC-01 (описание назначения документов), DOC-02 (настоящее ТЗ), DOC-03 (ПУКПО — план управления конфигурацией), DOC-04 (процедура формальной инспекции), DOC-05 (системные требования), DOC-06 и DOC-06.1--06.6 (программные требования — сводно и по подсистемам), DOC-07 (архитектура ПО), DOC-08 (справочник REST API), DOC-09 (матрица трассировки SR~$\to$~SWR~$\to$~TC), DOC-10 (план верификации ПО), DOC-11 (план интеграционного тестирования), DOC-13.1--13.6 (протоколы формальных инспекций), DOC-14 (количественный анализ покрытия требований).

\paragraph{\texttt{[TS.DOC.TRACE]}.} Требования должны быть трассируемы: каждое системное требование (SR-XXX) — минимум на одно программное требование (SWR-XXX) и минимум на один верификационный кейс (TC-XXX); каждое SWR — минимум на одно SR и один TC. Цепочка SR~$\to$~SWR~$\to$~TC ведётся в матрице трассировки DOC-09.

\paragraph{\texttt{[TS.DOC.TECH]}.} В качестве технологий реализации фиксируются: язык Kotlin 1.9.25, платформа Spring Boot 3.4.0 (Web, JPA, Security, Validation, Actuator), СУБД PostgreSQL 17, аутентификация JWT (jjwt), развёртывание Docker/docker-compose.

\paragraph{\texttt{[TS.DOC.NOTEST]}.} Тестовый код не пишется и прогоны не выполняются (решение от 20.10.2026): верификация обеспечивается формальными инспекциями и описательными планами DOC-10/DOC-11 без исполнения тестовых прогонов.

\section{Порядок контроля и приёмки \texttt{[TS.TEST]}}
Контроль и приёмка осуществляются на основании комплекса верификационных мероприятий по документации и Системе; исполнение тестовых прогонов не предусматривается (п.~4).

\subsection{Мероприятия по верификации \texttt{[TS.TEST.VER]}}
\begin{description}[style=nextline]
    \item[Формальные инспекции:] инспекции документов комплекта по процедуре DOC-04 (по Фагану), с ротацией рецензентов по чужим доменам; по документам приоритета оформляются протоколы DOC-13.1--13.6; выявленные несоответствия получают идентификаторы DEF-XXX и закрываются до конца спринта, следующего за протоколом.
    \item[Анализ трассировки:] проверка полноты цепочки SR~$\to$~SWR~$\to$~TC в матрице DOC-09; нарушение инвариантов трассировки (п.~4, \texttt{[TS.DOC.TRACE]}) является блокирующим замечанием.
    \item[Верификационные кейсы:] план верификации DOC-10 определяет описательные кейсы TC-XXX — позитивные и негативные (401 без JWT, 403 без роли, 400 при невалидном теле) — с ожидаемыми результатами по каждому требованию; исполнение прогонов не выполняется.
    \item[Сценарии интеграционного тестирования:] план DOC-11 описывает сквозные цепочки (регистрация $\to$ группа $\to$ сессия $\to$ решение $\to$ вердикт; автопровал по таймеру; редакционные генерации; модерация) со ссылками на кейсы TC-XXX из DOC-10, без исполнения прогонов.
\end{description}

\subsection{Критерии приёмки \texttt{[TS.TEST.ACCEPT]}}
Работа считается принятой при выполнении следующих условий:
\begin{itemize}
    \item все документы комплекта доведены до версии v1.0 и утверждены (листы регистрации изменений заполнены);
    \item формальные инспекции по всем документам приоритета проведены с результатом ПРИНЯТО или ПРИНЯТО С ЗАМЕЧАНИЯМИ (протоколы DOC-13.1--13.6 оформлены);
    \item все замечания DEF-XXX закрыты в установленные сроки;
    \item матрица трассировки DOC-09 полна — инварианты «каждое SR имеет минимум одно SWR и один TC; каждое SWR — минимум одно SR и один TC» выполнены;
    \item инварианты комплекта, установленные Реестром документов, выполнены; полнота покрытия подтверждена количественным анализом DOC-14.
\end{itemize}

\end{document}
